\scanpage{13}
\mspage{(11)}

plus précisément appeler \uline{complexe de De Rham à p.d. virtuel} \inserted{\uline{sur $k$}} \deleted{\unsure{un} complexe de} $k$-bi-algèbre diff. \unsure{augmentée}, ass. unit. alternée, à diff de bidegré $-1,+1$ \deleted{(ou $0,+1$ au choix)} à bidegrés $\geq 0$, avec p.d. sur l'idéal d'augm., compatibles avec la diff $(d(x^{(n)})=x^{(n-1)}dx)$, et telle que les $H^{\bullet,q}(C^{\bullet\bullet})[-q]$ soient des $q$-ombres\unsure{?} (auquel cas on récupère à partir de $C^{\bullet\bullet}$ un élément de $D(k)$ avec structure multiplicative ass. unit. comm.\dots), passer à une ``catégorie dérivée'' de ces complexes en ``inversant'' les flèches qui sont des \unsure{qu-isom}, d'où par $C^{\bullet\bullet}_{\mathrm{DRpd}}(\mathcal X_\bullet,k)$ un foncteur
\diagram{\begin{tikzcd}[column sep=large]
\begin{array}{c}(\mathrm{Hot})\\ \text{types d'hom}\end{array}
\arrow[r] &
\begin{array}{c}(\mathrm{DRpd})\\ \text{cat. dérivée}\\ \text{des complexes}\\ \mathrm{DRpd}\end{array}
\end{tikzcd}}
et on peut se demander si sa restriction aux espaces connexes et simplement connexes avec des $H^i_{\unsure{*}}(X,\mathbb Z)$ de t.f. (cas $k=\mathbb Z$) \deleted{est} \inserted{induit} une équivalence avec les complexes de DRàpd sur $\mathbb Z$ tels que $H^0(C^{\bullet\bullet})\simeq k$, $H^1(C^{\bullet\bullet})=0$, les $H^i(C^{\bullet\bullet})$ de t.f. sur $\mathbb Z$\dots

\mspage{12}
